% discussion.tex (국문 v10). 소제목 없음(Sci Rep).
\section{고찰}\label{sec:discussion}
이 연구는 라벨, 기후 입력, 채점이 같은 셀 · 같은 해를 가리키는 셀-연도 설계로 ALT 모델을 평가하였다. 결과는 세 가지로 정리된다. 첫째, 관측 연도의 기후를 관측일까지 누적해 입력하면, Stefan 잔차를 학습하는 gradient boosting 모델이 Stefan 모델의 오차를 레나 델타에서 \ndrefLeblkRoabs{} cm, 캐나다에서 \ndrefCablkRoabs{} cm 줄였고 알래스카에서는 같은 수준이었다. 둘째, 같은 모델은 한 지역을 학습에서 빼거나 라벨이 몇 개뿐일 때도 안정적이었고, 직접 ML은 학습에서 뺀 레나 델타에서 실패하였으며 더 많은 라벨이 필요하였다. 셋째, 같은 셀 · 같은 해 비교에서 이 모델은 레나 델타에서 공개 제품 5종보다 정확하였고, 알래스카와 캐나다에서는 두 ML 제품과 같은 수준이었다.

물리 잔차 구조가 작동하는 이유는 두 부분이 서로 다른 정보를 맡기 때문이다. Stefan 항은 해마다의 융해기 기후에 대한 반응을 맡고, 물리적으로 제한된 형태로 새 지역과 새 해에 외삽된다. 잔차 모델은 토양과 지형 때문에 생기는 지역 계수의 편차를 보정한다. 이득에는 지형과 토양 공변량이 필요하였고, Stefan 오차가 가장 큰 조건에서 가장 컸으며, 잔차 모델의 기여 분석은 레나 델타에서 지형, 캐나다에서 토양 유기탄소를 가리켰다. Stefan 유사라벨을 더하면 학습기가 Stefan 모델 쪽으로 끌려가 정확도가 늘지 않았고, 기후 반응을 라벨에서 직접 배워야 하는 직접 ML은 학습에서 뺀 레나 델타에서 정확도를 \ndrefLeregDzabs{} cm 잃었다. 한 지역의 라벨을 모두 쓰면 직접 ML도 잔차 모델과 비슷하거나 낮은 점추정 오차에 이를 수 있다. 잔차 구조의 장점은 그 정확도가 같은 지역의 많은 라벨에 의존하지 않는다는 점이다.

기존 제품이 보고한 정확도는 서로 다른 지점 집합과 검증 방식에 따른 값이어서(예: Liu 등의 무작위 10겹 RMSE 21.6 cm\cite{Liu2024}, Wei 등의 지점 하나 제외 RMSE 60.5 cm\cite{Wei2026}) 서로 비교할 수 없다. 모델과 제품을 같은 셀 · 같은 해에서 채점하면 이 모호함이 없어진다. 이 비교에서 나타난 지역 차이는 제품의 학습 자료와 대응한다. Wei 등과 Liu 등은 알래스카 라벨 셀에서 몇 km 안에 있는 CALM, GTN-P, ABoVE 관측으로 학습하였고, 알래스카와 캐나다에서 우리 모델과 같은 정확도를 보였다. 레나 델타는 Wei 등의 가장 가까운 학습 지점에서 약 \ndistweiLemed{} km 떨어져 있고, 두 제품 모두 이곳의 얇은 활동층을 과대 추정한다. 레나 델타 라벨 없이 학습한 모델도 이곳에서 이미 Wei 등보다 정확하였고, 레나 델타 라벨을 쓰면 오차가 \nccLeRoall{} cm로 낮아져 제품의 \nccLeCcifiveall--\nccLeWeiall{} cm보다 작았다. 잔차 학습기는 제품을 기준값으로 쓸 때에도 정확도를 높였다. 따라서 지역 라벨과 잔차 학습을 결합하면 Stefan 모델뿐 아니라 기존 제품도 그 제품이 학습하지 않은 지역에 맞춰 보정할 수 있다.

라벨 실험은 새 지역에 적용하는 순서를 제시한다. 지역 라벨이 없을 때는 직접 ML이 실패할 수 있으므로 다른 지역 계수의 Stefan 모델이 더 안전한 출발점이다. 라벨이 25개 정도 모이면 Stefan 계수를 다시 맞추고 잔차를 학습해 레나 델타의 정확도를 회복하였다. 외부 CALM 지역에서는 라벨 5개로 오차가 \nftAllRszero{} cm에서 \nftAllRsfive{} cm로 낮아져 ESA CCI v5 제품보다 작아졌다. 외부 지점은 전이의 한계도 보여 준다. 서 · 중부 시베리아와 스발바르의 융해 깊이는 연구 지역의 범위를 넘고, 그곳의 이득은 잔차 모델이 아니라 계수 재보정에서 나왔다. 캐나다는 반대의 경우로, 원천 지역으로 학습한 직접 ML이 더 정확하였고 적은 캐나다 표본으로 다시 맞춘 계수는 오차를 키웠다. 따라서 필요한 지역 라벨 수는 새 지역이 학습 범위에서 얼마나 떨어져 있는지에 따라 달라지며, 연구 지역과 비슷한 지역에서는 토양과 지형 범위를 덮는 라벨 25--100개가 적절한 목표이다.

연도가 명시된 지도는 이 설계의 두 번째 산출물이다. 각 지도는 명시된 해의 ERA5-Land 기후로 계산되므로 같은 해의 관측이나 다른 제품과 비교할 수 있고, 편차를 해마다 추적할 수 있다(예: 알래스카 \nanmaxyrAk{}년 \nanmaxAk{} cm). 지도는 적용 가능 영역과 90\% 예측 구간을 함께 제공하여 예측이 학습 조건을 벗어나는 셀을 표시한다. 예측의 연도 간 범위는 관측 셀 평균의 범위보다 작다(장기 관측 셀은 보충 그림 S10). 관측 셀이 캠페인마다 바뀌고, 잔차 모델이 주로 공간 대비로 학습되기 때문이다.

한계는 다음과 같다. 알래스카 라벨의 절반이 2014년 자료여서 그곳에서는 관측 연도 입력의 효과를 판정하지 못하였다. 레나 델타 라벨은 한 자료원의 \nnblkLe{}개 블록에서 나와 구간이 넓다. 기후 입력은 월자료이고 토양 공변량은 정적이어서 짧은 기상 사건과 토양 수분 변화는 반영되지 않는다. 주로 산지와 영구동토 지대의 남쪽 가장자리에 있는 약 3 m보다 두꺼운 활동층은 학습 범위 밖이다. 일 단위 기후 입력, 적설 · 토양 수분 정보, 그리고 CALM과 집중 관측 자료가 함께 있는 서시베리아 같은 네 번째 지역의 라벨이 다음 단계이다.

요약하면, 지역 라벨로 재보정하고 잔차 학습기로 보정한 Stefan 모델은 토양과 지형이 중요한 곳에서 Stefan 모델보다 정확하고, 라벨이 없는 지역에서도 안정적이며, 적은 라벨로 적응하는 연도별 ALT를 제공한다. 같은 셀 · 같은 해로 평가하면 이 모델은 제품이 학습하지 않은 지역에서 공개 제품보다 정확하고, 제품이 학습한 지역에서는 같은 수준이다. 같은 절차는 다른 영구동토 지역에 적용할 수 있고, 기존 제품을 기준값으로 쓸 수도 있으며, 적용 가능 영역이 명시된 연도별 1 km 지도로 활동층 감시에 쓰일 수 있다.
